\chapter{§3.4 三角形}
\label{sec:3.4}


\prerequisite{§3.2, §3.3}

\gradelevel{7-8 年级}


\section*{三角形的基本概念}


\begin{explore}


自行车的车架是三角形的——为什么不做成四边形？试着用三根木条钉成一个三角形框架，再用四根木条钉成四边形框架，分别推一推。你会发现三角形推不动，而四边形一推就变形了。这就是三角形的\textbf{稳定性}。

\end{explore}

\begin{understand}


由三条线段首尾依次相连围成的图形叫做\textbf{三角形}。三角形 $ABC$ 记作 $\triangle ABC$ 。

三角形的基本要素：
\begin{itemize}
  \item \textbf{顶点}： $A$ 、 $B$ 、 $C$
  \item \textbf{边}： $AB$ 、 $BC$ 、 $AC$ （也常用小写字母 $a$ 、 $b$ 、 $c$ 表示对边，其中 $a = BC$ 、 $b = AC$ 、 $c = AB$ ）
  \item \textbf{角}： $\angle A$ 、 $\angle B$ 、 $\angle C$ （即 $\angle BAC$ 、 $\angle ABC$ 、 $\angle ACB$ ）
\end{itemize}


\textbf{三角形的分类}

按角分类：
\begin{center}
\begin{tabular}{ll}
\toprule
类型 & 条件 \\
\midrule
锐角三角形 & 三个角都是锐角 \\
直角三角形 & 有一个角是直角 \\
钝角三角形 & 有一个角是钝角 \\
\bottomrule
\end{tabular}
\end{center}


按边分类：
\begin{center}
\begin{tabular}{ll}
\toprule
类型 & 条件 \\
\midrule
不等边三角形 & 三条边都不相等 \\
等腰三角形 & 至少有两条边相等 \\
等边三角形 & 三条边都相等 \\
\bottomrule
\end{tabular}
\end{center}


\textbf{三角形的内角和}

$\triangle ABC$ 的三个内角之和等于 $180°$ ：

\[
\angle A + \angle B + \angle C = 180°
\]


推理过程（参见图 assets/03-geometry/triangle-angle-sum.svg）：

过顶点 $A$ 作直线 $DE \parallel BC$ 。

因为 $DE \parallel BC$ ，所以 $\angle DAB = \angle B$ （内错角相等）， $\angle EAC = \angle C$ （内错角相等）。

因为 $\angle DAB + \angle BAC + \angle EAC = 180°$ （平角），

所以 $\angle B + \angle A + \angle C = 180°$ 。

\begin{center}
\begin{tikzpicture}[every node/.style={font=\small}]
  \coordinate (A) at (2.5, 2.8);
  \coordinate (B) at (0,   0);
  \coordinate (C) at (5,   0);
  \draw[main line] (A) -- (B) -- (C) -- cycle;
  \node[above] at (A) {$A$};
  \node[left]  at (B) {$B$};
  \node[right] at (C) {$C$};

  % Parallel line through A (dashed)
  \draw[ext line] (-0.5,2.8) -- (5.5,2.8);

  % Angle alpha at B (color: red)
  \draw[angle arc, draw=red!70!black, thick]
    ($(B)+(20:0.6)$) arc[start angle=20, end angle=62, radius=0.6];
  \node[red!70!black] at ($(B)+(41:0.95)$) {$\alpha$};

  % Angle gamma at A vertex
  \draw[angle arc, draw=green!60!black, thick]
    ($(A)+(-62:0.5)$) arc[start angle=-62, end angle=-118, radius=0.5];
  \node[green!60!black] at ($(A)+(270:0.8)$) {$\gamma$};

  % Angle beta at C (color: blue)
  \draw[angle arc, draw=blue!70!black, thick]
    ($(C)+(118:0.6)$) arc[start angle=118, end angle=160, radius=0.6];
  \node[blue!70!black] at ($(C)+(139:0.95)$) {$\beta$};

  % alpha' on parallel line left of A (= alpha, alternate interior)
  \draw[angle arc, draw=red!70!black, thick]
    ($(A)+(-118:0.5)$) arc[start angle=-118, end angle=-180, radius=0.5];
  \node[red!70!black] at ($(A)+(-149:0.85)$) {$\alpha$};

  % beta' on parallel line right of A (= beta)
  \draw[angle arc, draw=blue!70!black, thick]
    ($(A)+(0:0.5)$) arc[start angle=0, end angle=-62, radius=0.5];
  \node[blue!70!black] at ($(A)+(-31:0.85)$) {$\beta$};

  \node[below] at (2.5,-0.6) {$\alpha + \gamma + \beta = 180°$};
\end{tikzpicture}
\captionof{figure}{三角形内角和定理证明：过 $A$ 作 $BC$ 的平行线}
\label{fig:triangle-angle-sum}
\end{center}


\end{understand}

\begin{workedexamples}


\textbf{例 1.} $\triangle ABC$ 中， $\angle A = 50°$ ， $\angle B = 70°$ ，求 $\angle C$ 。

\textbf{解：}
$\angle C = 180° - \angle A - \angle B = 180° - 50° - 70° = 60°$ 。

\textbf{例 2.} $\triangle ABC$ 中， $\angle A : \angle B : \angle C = 2 : 3 : 4$ ，求三个角。

\textbf{解：}
设 $\angle A = 2k$ ， $\angle B = 3k$ ， $\angle C = 4k$ 。

$2k + 3k + 4k = 180°$ ， $9k = 180°$ ， $k = 20°$ 。

所以 $\angle A = 40°$ ， $\angle B = 60°$ ， $\angle C = 80°$ 。

\textbf{例 3.} 直角三角形中，一个锐角为 $35°$ ，求另一个锐角。

\textbf{解：}
设直角为 $\angle C = 90°$ ，已知锐角 $\angle A = 35°$ 。

$\angle B = 180° - 90° - 35° = 55°$ 。

\textbf{例 4.} $\triangle ABC$ 中， $\angle A = 2\angle B = 3\angle C$ ，求三个角。

\textbf{解：}
设 $\angle C = x$ ，则 $\angle A = 3x$ ，由 $\angle A = 2\angle B$ 得 $\angle B = \dfrac{3x}{2}$ 。

$3x + \dfrac{3x}{2} + x = 180°$

$\dfrac{6x + 3x + 2x}{2} = 180°$

$\dfrac{11x}{2} = 180°$ ， $x = \dfrac{360°}{11} \approx 32.7°$ 。

精确值： $\angle C = \dfrac{360°}{11}$ ， $\angle B = \dfrac{540°}{11}$ ， $\angle A = \dfrac{1080°}{11}$ 。

\end{workedexamples}

\begin{keytakeaway}


\begin{itemize}
  \item 三角形的内角和为 $180°$ 。
  \item 直角三角形的两个锐角互余。
  \item 三角形具有稳定性。
\end{itemize}


\end{keytakeaway}

\begin{practice}


\begin{enumerate}
  \item $\triangle ABC$ 中， $\angle A = 48°$ ， $\angle C = 65°$ ，求 $\angle B$ 。
  \item 一个等腰三角形的顶角是 $40°$ ，求两个底角。
\end{enumerate}


\end{practice}



\section*{三角形的三边关系}


\begin{explore}


给你三根木棒，长度分别是 $3$ cm、 $4$ cm、 $5$ cm，能搭成三角形吗？如果换成 $1$ cm、 $2$ cm、 $5$ cm呢？试一试就知道——第二组搭不成！两条短边加起来还不够长，无法"够到"第三条边的两端。

\end{explore}

\begin{understand}


\textbf{三角形的三边关系}：三角形任意两边之和大于第三边。

\[
a + b > c, \quad a + c > b, \quad b + c > a
\]


等价地，也可以说：三角形任意两边之差小于第三边。

\end{understand}

\begin{quote}
实际判断时，只需检验\textbf{两条较短边之和是否大于最长边}即可。如果最短的两边之和都大于最长边，那么其他两个不等式自然成立。
\end{quote}


\begin{workedexamples}


\textbf{例 1.} 判断下列各组线段能否构成三角形：(a) $3, 4, 8$ ；(b) $5, 6, 10$ ；(c) $3, 3, 5$ 。

\textbf{解：}
(a) $3 + 4 = 7 < 8$ ，不能构成三角形。

(b) $5 + 6 = 11 > 10$ ，能构成三角形。

(c) $3 + 3 = 6 > 5$ ，能构成三角形。

\textbf{例 2.} 三角形的两边长分别为 $3$ cm 和 $7$ cm，第三边的长度 $x$ 的取值范围是什么？

\textbf{解：}
由三边关系： $7 - 3 < x < 7 + 3$ ，即 $4 < x < 10$ 。

\textbf{例 3.} 等腰三角形的两边长分别为 $3$ 和 $7$ ，求周长。

\textbf{解：}
等腰三角形有两条边相等。

情况一：两腰为 $3$ ，底边为 $7$ 。检验： $3 + 3 = 6 < 7$ ，不满足三边关系，不能构成三角形。

情况二：两腰为 $7$ ，底边为 $3$ 。检验： $7 + 3 = 10 > 7$ ，满足三边关系。

所以周长 $= 7 + 7 + 3 = 17$ 。

\end{workedexamples}

\begin{keytakeaway}


\begin{itemize}
  \item 三角形任意两边之和大于第三边。
  \item 判断时只需检验两短边之和 $>$ 最长边。
  \item 已知两边求第三边范围：两边之差 $<$ 第三边 $<$ 两边之和。
\end{itemize}


\end{keytakeaway}

\begin{practice}


\begin{enumerate}
  \item 三角形的两边长分别为 $5$ 和 $9$ ，第三边为整数，求第三边的所有可能值。
  \item 等腰三角形的两边长为 $4$ 和 $9$ ，求周长。
\end{enumerate}


\end{practice}



\section*{全等三角形}


\begin{explore}


冲压模具生产出的零件，每一个的形状和大小都完全一样——它们是"全等"的。两个三角形，如果能完全重合，我们就说它们全等。

\end{explore}

\begin{understand}


如果两个三角形的形状和大小完全相同（能够完全重合），则称这两个三角形\textbf{全等}，记作 $\triangle ABC \cong \triangle DEF$ 。

全等三角形的对应关系：
\begin{itemize}
  \item \textbf{对应顶点}： $A \leftrightarrow D$ ， $B \leftrightarrow E$ ， $C \leftrightarrow F$
  \item \textbf{对应边相等}： $AB = DE$ ， $BC = EF$ ， $AC = DF$
  \item \textbf{对应角相等}： $\angle A = \angle D$ ， $\angle B = \angle E$ ， $\angle C = \angle F$
\end{itemize}


\end{understand}

\begin{quote}
书写全等时，对应顶点要写在对应位置上。 $\triangle ABC \cong \triangle DEF$ 说明 $A$ 对应 $D$ 、 $B$ 对应 $E$ 、 $C$ 对应 $F$ 。
\end{quote}


\textbf{全等三角形的判定条件}

要判断两个三角形全等，不需要知道全部六个元素（三边三角），只需要满足以下条件之一：

\begin{center}
\begin{tabular}{lll}
\toprule
判定方法 & 条件 & 简记 \\
\midrule
边边边 & 三条边分别相等 & SSS \\
边角边 & 两边及其夹角分别相等 & SAS \\
角边角 & 两角及其夹边分别相等 & ASA \\
角角边 & 两角及其中一角的对边分别相等 & AAS \\
斜边直角边 & 直角三角形的斜边和一条直角边分别相等 & HL \\
\bottomrule
\end{tabular}
\end{center}


\begin{quote}\textbf{注意：}\textbf{SSA（边边角）不能判定全等！} 已知两边和非夹角相等时，可能存在两种不同的三角形。这是一个常见的错误，务必记住。\end{quote}


\begin{center}
\begin{tikzpicture}[every node/.style={font=\small}, scale=0.9]
  % Triangle 1: ABC with AB=4, BC=3, angle A = 45°
  \coordinate (A1) at (0,0);
  \coordinate (B1) at (4,0);
  % angle A = 45°, AB=4, BC=3: two possible positions for C
  % C1 above (larger triangle)
  \coordinate (C1) at ($(A1)+(45:5.66)$);
  % Find C such that BC = 3: C on ray from A at 45° intersecting circle of radius 3 from B
  % Actually use concrete coords: C = (2.12, 2.12) approx; then |BC|=sqrt((4-2.12)^2+2.12^2)=sqrt(3.54+4.49)=sqrt(8.03)≈2.83... not 3
  % Use: A=(0,0), angle BAC=45°, AB=4. Ray from A: direction (cos45°,sin45°)=(0.707,0.707)
  % On this ray: P(t) = (0.707t, 0.707t). Circle center B=(4,0), r=3.
  % |P-B|^2 = (0.707t-4)^2+(0.707t)^2 = 3^2=9
  % 0.5t^2 - 5.657t + 16 + 0.5t^2 = 9 → t^2 - 5.657t + 7 = 0
  % t = (5.657 ± sqrt(32-28))/2 = (5.657 ± 2)/2 → t1=3.828, t2=1.828
  \coordinate (C1) at (2.707, 2.707);   % t=3.828: (3.828*0.707, 3.828*0.707)
  \coordinate (C2) at (1.293, 1.293);   % t=1.828: (1.828*0.707, 1.828*0.707)
  % Left triangle (C1)
  \draw[main line] (A1)--(B1)--(C1)--cycle;
  \node[below left] at (A1) {$A$};
  \node[below right] at (B1) {$B$};
  \node[above] at (C1) {$C_1$};
  % Right triangle (C2)
  \begin{scope}[xshift=7cm]
    \coordinate (A2) at (0,0);
    \coordinate (B2) at (4,0);
    \coordinate (C2b) at (1.293, 1.293);
    \draw[main line] (A2)--(B2)--(C2b)--cycle;
    \node[below left] at (A2) {$A$};
    \node[below right] at (B2) {$B$};
    \node[above left] at (C2b) {$C_2$};
  \end{scope}
  \node[below] at (3.5,-0.8) {$AB=AB,\ BC_1=BC_2,\ \angle A=\angle A$，但两三角形不全等！（SSA 无效）};
\end{tikzpicture}
\captionof{figure}{SSA 反例：相同的两边和非夹角无法唯一确定三角形}
\label{fig:ssa-counterexample}
\end{center}


\begin{workedexamples}


\textbf{例 1.} 如图， $AB = DC$ ， $\angle ABD = \angle DCB$ 。证明 $\triangle ABD \cong \triangle DCB$ 。

\begin{center}
\begin{tikzpicture}[every node/.style={font=\small}, scale=0.9]
  % Triangle ABC
  \coordinate (A) at (0,0);
  \coordinate (B) at (3,0);
  \coordinate (C) at (1,2.2);
  \draw[main line] (A)--(B)--(C)--cycle;
  \node[below left] at (A) {$A$};
  \node[below right] at (B) {$B$};
  \node[above] at (C) {$C$};
  % Single tick on AB
  \draw ($(A)!0.5!(B)$)++(90:0.1)--++(0,-0.2);
  % Double tick on AC
  \draw ($(A)!0.45!(C)$)++({90+atan2(2.2,1)}:0.1)--++({-90+atan2(2.2,1)}:0.1);
  \draw ($(A)!0.55!(C)$)++({90+atan2(2.2,1)}:0.1)--++({-90+atan2(2.2,1)}:0.1);
  % Angle A mark
  \draw[angle arc] ($(A)+(0:0.5)$) arc[start angle=0,end angle={atan2(2.2,1)},radius=0.5];

  % Triangle DEF (congruent, shifted right)
  \begin{scope}[xshift=5.5cm]
    \coordinate (D) at (0,0);
    \coordinate (E) at (3,0);
    \coordinate (F) at (1,2.2);
    \draw[main line] (D)--(E)--(F)--cycle;
    \node[below left] at (D) {$D$};
    \node[below right] at (E) {$E$};
    \node[above] at (F) {$F$};
    \draw ($(D)!0.5!(E)$)++(90:0.1)--++(0,-0.2);
    \draw ($(D)!0.45!(F)$)++({90+atan2(2.2,1)}:0.1)--++({-90+atan2(2.2,1)}:0.1);
    \draw ($(D)!0.55!(F)$)++({90+atan2(2.2,1)}:0.1)--++({-90+atan2(2.2,1)}:0.1);
    \draw[angle arc] ($(D)+(0:0.5)$) arc[start angle=0,end angle={atan2(2.2,1)},radius=0.5];
  \end{scope}

  \node at (4.7,1) {$\cong$};
\end{tikzpicture}
\captionof{figure}{SAS：$AB=DE$，$\angle A=\angle D$，$AC=DF$ $\Rightarrow$ $\triangle ABC \cong \triangle DEF$}
\label{fig:congruent-ex1}
\end{center}


\textbf{证明：}
在 $\triangle ABD$ 和 $\triangle DCB$ 中：

$$\begin{aligned}
AB &= DC \quad (\text{已知}) \\
\angle ABD &= \angle DCB \quad (\text{已知}) \\
BD &= CB \quad (\text{公共边})
\end{aligned}$$

所以 $\triangle ABD \cong \triangle DCB$ （SAS）。

\textbf{例 2.} 如图， $\angle A = \angle D$ ， $\angle B = \angle E$ ， $BC = EF$ 。证明 $\triangle ABC \cong \triangle DEF$ 。

\textbf{证明：}
在 $\triangle ABC$ 和 $\triangle DEF$ 中：

$$\begin{aligned}
\angle B &= \angle E \quad (\text{已知}) \\
BC &= EF \quad (\text{已知}) \\
\angle A &= \angle D \quad (\text{已知})
\end{aligned}$$

因为 $\angle B = \angle E$ ， $BC = EF$ ，又 $\angle A = \angle D$ ，

这里 $BC$ 是 $\angle B$ 的一条边， $\angle A$ 是 $\triangle ABC$ 的另一个角，

所以 $\triangle ABC \cong \triangle DEF$ （AAS）。

\textbf{例 3.} 如图， $\angle ACB = \angle ADB = 90°$ ， $AC = AD$ 。证明 $\triangle ACB \cong \triangle ADB$ 。

\begin{center}
\begin{tikzpicture}[every node/.style={font=\small}, scale=0.9]
  % Triangle ABC
  \coordinate (A) at (0,0);
  \coordinate (B) at (3.5,0);
  \coordinate (C) at (2,2.5);
  \draw[main line] (A)--(B)--(C)--cycle;
  \node[below left] at (A) {$A$};
  \node[below right] at (B) {$B$};
  \node[above] at (C) {$C$};
  % Mark: angle A (arc)
  \draw[angle arc] ($(A)+(0:0.55)$) arc[start angle=0,end angle={atan2(2.5,2)},radius=0.55];
  % Mark: angle B (arc)
  \draw[angle arc] ($(B)+(180:0.55)$) arc[start angle=180,end angle={180-atan2(2.5,1.5)},radius=0.55];
  % Single tick on BC (non-included side)
  \draw ($(B)!0.5!(C)$)++({90+atan2(2.5,-1.5)}:0.12)--++({-90+atan2(2.5,-1.5)}:0.12);

  % Triangle DEF
  \begin{scope}[xshift=5.5cm]
    \coordinate (D) at (0,0);
    \coordinate (E) at (3.5,0);
    \coordinate (F) at (2,2.5);
    \draw[main line] (D)--(E)--(F)--cycle;
    \node[below left] at (D) {$D$};
    \node[below right] at (E) {$E$};
    \node[above] at (F) {$F$};
    \draw[angle arc] ($(D)+(0:0.55)$) arc[start angle=0,end angle={atan2(2.5,2)},radius=0.55];
    \draw[angle arc] ($(E)+(180:0.55)$) arc[start angle=180,end angle={180-atan2(2.5,1.5)},radius=0.55];
    \draw ($(E)!0.5!(F)$)++({90+atan2(2.5,-1.5)}:0.12)--++({-90+atan2(2.5,-1.5)}:0.12);
  \end{scope}

  \node at (4.7,1.2) {$\cong$};
\end{tikzpicture}
\captionof{figure}{AAS：$\angle A=\angle D$，$\angle B=\angle E$，$BC=EF$ $\Rightarrow$ $\triangle ABC \cong \triangle DEF$}
\label{fig:congruent-ex3}
\end{center}


\textbf{证明：}
在直角 $\triangle ACB$ 和直角 $\triangle ADB$ 中：

$$\begin{aligned}
\angle ACB &= \angle ADB = 90° \quad (\text{已知}) \\
AB &= AB \quad (\text{公共斜边}) \\
AC &= AD \quad (\text{已知})
\end{aligned}$$

所以 $\triangle ACB \cong \triangle ADB$ （HL）。

\end{workedexamples}

\begin{keytakeaway}


\begin{itemize}
  \item 全等三角形对应边相等，对应角相等。
  \item 五种判定方法：SSS、SAS、ASA、AAS、HL。
  \item SSA 不能判定全等。
  \item 写全等时对应顶点写在对应位置。
\end{itemize}


\end{keytakeaway}

\begin{practice}


\begin{enumerate}
  \item 如图， $AB = CD$ ， $AB \parallel CD$ 。证明 $\triangle ABO \cong \triangle CDO$ （ $O$ 是 $AC$ 与 $BD$ 的交点）。
  \item 如图， $\angle 1 = \angle 2$ ， $\angle 3 = \angle 4$ ， $AD = BC$ 。证明 $\triangle ADB \cong \triangle BCA$ 。
  \item 如图， $RT\triangle ABC$ 和 $RT\triangle DEF$ 中， $\angle C = \angle F = 90°$ ， $AB = DE$ ， $BC = EF$ 。证明 $\triangle ABC \cong \triangle DEF$ 。
\end{enumerate}


\end{practice}



\section*{等腰三角形与等边三角形}


\begin{explore}


用一张纸对折后沿斜线剪一刀，展开后得到的三角形两边相等——这就是等腰三角形。如果剪切角度恰当，还能得到三边都相等的等边三角形。这些特殊三角形有哪些特别的性质？

\end{explore}

\begin{understand}


\textbf{等腰三角形}：有两条边相等的三角形。相等的两条边叫做\textbf{腰}，第三条边叫做\textbf{底边}，两腰所夹的角叫做\textbf{顶角}，底边上的两个角叫做\textbf{底角}。

\begin{center}
\begin{tikzpicture}[every node/.style={font=\small}]
  \coordinate (A) at (2, 3.5);
  \coordinate (B) at (0, 0);
  \coordinate (C) at (4, 0);
  \draw[main line] (A)--(B)--(C)--cycle;
  \node[above] at (A) {$A$};
  \node[below left]  at (B) {$B$};
  \node[below right] at (C) {$C$};
  % Tick marks: AB = AC (double tick)
  \draw ($(A)!0.45!(B)$)++({90+atan2(3.5,2)}:0.12)--++({-90+atan2(3.5,2)}:0.12);
  \draw ($(A)!0.55!(B)$)++({90+atan2(3.5,2)}:0.12)--++({-90+atan2(3.5,2)}:0.12);
  \draw ($(A)!0.45!(C)$)++({90-atan2(3.5,2)}:0.12)--++({-90+atan2(3.5,2)}:0.12);
  \draw ($(A)!0.55!(C)$)++({90-atan2(3.5,2)}:0.12)--++({-90+atan2(3.5,2)}:0.12);
  % Base angles B = C
  \draw[angle arc, draw=red!70!black]
    ($(B)+({atan2(3.5,2)}:0.6)$) arc[start angle={atan2(3.5,2)},end angle=0,radius=0.6];
  \draw[angle arc, draw=red!70!black]
    ($(C)+({180-atan2(3.5,2)}:0.6)$) arc[start angle={180-atan2(3.5,2)},end angle=180,radius=0.6];
  \node[red!70!black, font=\scriptsize] at ($(B)+(20:1.0)$) {$\angle B$};
  \node[red!70!black, font=\scriptsize] at ($(C)+(160:1.0)$) {$\angle C$};
\end{tikzpicture}
\captionof{figure}{等腰三角形：$AB=AC$，底角 $\angle B = \angle C$（等腰三角形的底角相等）}
\label{fig:isosceles-triangle}
\end{center}


\textbf{等腰三角形的性质}：

\begin{enumerate}
  \item \textbf{等边对等角}（等腰三角形的两个底角相等）。
  \item \textbf{"三线合一"}：等腰三角形的顶角平分线、底边上的中线、底边上的高互相重合。
\end{enumerate}


\textbf{等腰三角形的判定}：

\textbf{等角对等边}：如果一个三角形有两个角相等，则这两个角所对的边也相等（该三角形是等腰三角形）。

\textbf{等边三角形}：三条边都相等的三角形。等边三角形是等腰三角形的特殊情况。

等边三角形的性质：
\begin{itemize}
  \item 三个角都等于 $60°$ 。
  \item 等边三角形也满足"三线合一"。
\end{itemize}


等边三角形的判定：
\begin{itemize}
  \item 三边相等的三角形是等边三角形。
  \item 三个角都相等的三角形是等边三角形。
  \item 有一个角是 $60°$ 的等腰三角形是等边三角形。
\end{itemize}


\end{understand}

\begin{workedexamples}


\textbf{例 1.} 等腰三角形的一个角是 $80°$ ，求其他两个角。

\textbf{解：}
情况一： $80°$ 是顶角。两底角 $= \dfrac{180° - 80°}{2} = 50°$ 。

情况二： $80°$ 是底角。另一底角也是 $80°$ ，顶角 $= 180° - 80° - 80° = 20°$ 。

所以其他两个角为 $50°, 50°$ 或 $80°, 20°$ 。

\textbf{例 2.} 等腰三角形的一个角是 $120°$ ，求其他两个角。

\textbf{解：}
$120°$ 只能是顶角（因为底角相等，若两个底角都是 $120°$ ，角度之和已超过 $180°$ ）。

两底角 $= \dfrac{180° - 120°}{2} = 30°$ 。

\textbf{例 3.} 如图， $\triangle ABC$ 中， $AB = AC$ ， $BD$ 是 $\angle ABC$ 的角平分线， $\angle A = 40°$ 。求 $\angle BDC$ 。

\begin{center}
\begin{tikzpicture}[every node/.style={font=\small}]
  \coordinate (A) at (2.5, 3.5);
  \coordinate (B) at (0, 0);
  \coordinate (C) at (5, 0);
  \coordinate (M) at (2.5, 0);  % midpoint of BC
  \draw[main line] (A)--(B)--(C)--cycle;
  \draw[aux line] (A)--(M);
  \node[above] at (A) {$A$};
  \node[below left]  at (B) {$B$};
  \node[below right] at (C) {$C$};
  \node[below] at (M) {$M$};
  \fill (M) circle (2pt);
  % Right angle at M
  \draw (2.72,0)--(2.72,0.22)--(2.5,0.22);
  % Equal tick marks on BM and MC
  \draw (1.15,-0.12)--(1.35,0.12);
  \draw (3.65,-0.12)--(3.85,0.12);
  \node[right, font=\scriptsize] at (2.8,1.8) {中线 = 高 = 角平分线};
\end{tikzpicture}
\captionof{figure}{等腰三角形中，顶角的角平分线、底边中线、底边上的高三线合一}
\label{fig:isosceles-ex3}
\end{center}


\textbf{解：}
因为 $AB = AC$ ， $\angle A = 40°$ ，

所以 $\angle ABC = \angle ACB = \dfrac{180° - 40°}{2} = 70°$ 。

因为 $BD$ 平分 $\angle ABC$ ，

所以 $\angle ABD = \dfrac{70°}{2} = 35°$ 。

在 $\triangle BDC$ 中， $\angle BDC = 180° - \angle DBC - \angle BCD = 180° - 35° - 70° = 75°$ 。

\textbf{例 4.} 证明：等边三角形的每个角都是 $60°$ 。

\textbf{证明：}
设 $\triangle ABC$ 是等边三角形， $AB = BC = CA$ 。

因为 $AB = AC$ ，所以 $\angle B = \angle C$ （等边对等角）。

因为 $AB = BC$ ，所以 $\angle A = \angle C$ （等边对等角）。

所以 $\angle A = \angle B = \angle C$ 。

又因为 $\angle A + \angle B + \angle C = 180°$ ，所以 $3\angle A = 180°$ ， $\angle A = 60°$ 。

因此等边三角形的每个角都是 $60°$ 。

\end{workedexamples}

\begin{keytakeaway}


\begin{itemize}
  \item 等腰三角形：等边对等角，等角对等边。
  \item 等腰三角形有"三线合一"的性质。
  \item 等边三角形三个角都是 $60°$ 。
  \item 等腰三角形的角分两种情况讨论（给定角是顶角还是底角）。
\end{itemize}


\end{keytakeaway}

\begin{practice}


\begin{enumerate}
  \item 等腰三角形的一个角是 $50°$ ，求其他两个角（考虑所有情况）。
  \item $\triangle ABC$ 中， $AB = AC$ ， $D$ 是 $BC$ 的中点。证明 $AD \perp BC$ 。
  \item 已知等边 $\triangle ABC$ 的周长为 $18$ cm，求边长。
\end{enumerate}


\end{practice}



\section*{相似三角形}


\begin{explore}


把一张照片放大或缩小后，照片中的人物比例不变——长宽比是一样的。这就是"相似"的直观含义：形状相同，大小可以不同。

\end{explore}

\begin{understand}


如果两个三角形的对应角相等、对应边成比例，则这两个三角形\textbf{相似}，记作 $\triangle ABC \sim \triangle DEF$ 。

对应边的比值叫做\textbf{相似比}（或比例系数）。若 $\triangle ABC \sim \triangle DEF$ ，相似比为 $k$ ，则：

\[
\frac{AB}{DE} = \frac{BC}{EF} = \frac{AC}{DF} = k
\]


\end{understand}

\begin{quote}
全等是相似的特殊情况——相似比 $k = 1$ 时，就是全等。
\end{quote}


\textbf{相似三角形的判定条件}：

\begin{center}
\begin{tabular}{ll}
\toprule
判定方法 & 条件 \\
\midrule
AA & 两组对应角分别相等 \\
SAS 相似 & 两组对应边成比例，且夹角相等 \\
SSS 相似 & 三组对应边分别成比例 \\
\bottomrule
\end{tabular}
\end{center}


\begin{quote}
AA 判定最常用：只需两对角相等（第三对角自动相等，因为内角和为 $180°$ ）。
\end{quote}


\textbf{相似三角形的性质}：
\begin{itemize}
  \item 对应角相等。
  \item 对应边成比例。
  \item 对应高的比、对应中线的比、对应角平分线的比都等于相似比。
  \item 周长的比等于相似比。
  \item 面积的比等于相似比的平方。
\end{itemize}


\begin{workedexamples}


\textbf{例 1.} 如图， $DE \parallel BC$ ， $AD = 3$ ， $DB = 6$ ， $DE = 4$ 。求 $BC$ 的长。

\begin{center}
\begin{tikzpicture}[every node/.style={font=\small}]
  % Large triangle ABC: vertices at concrete positions (3-4-5 scaled ×2)
  \coordinate (A) at (1.2, 4);
  \coordinate (B) at (0, 0);
  \coordinate (C) at (3, 0);
  \draw[main line] (A)--(B)--(C)--cycle;
  \node[above] at (A) {$A$};
  \node[below left] at (B) {$B$};
  \node[below right] at (C) {$C$};
  % Side labels
  \node[left]  at ($(A)!0.5!(B)$) {$6$};
  \node[below] at ($(B)!0.5!(C)$) {$8$};
  \node[right] at ($(A)!0.5!(C)$) {$10$};
  % Angle arc at A (similar mark)
  \draw[angle arc] ($(A)+(-100:0.4)$) arc[start angle=-100,end angle=-145,radius=0.4];
  % Angle arc at B
  \draw[angle arc] ($(B)+(60:0.4)$) arc[start angle=60,end angle=0,radius=0.4];

  % Small triangle DEF (half scale, shifted right)
  \begin{scope}[xshift=5.5cm]
    \coordinate (D) at (0.6, 2);
    \coordinate (E) at (0, 0);
    \coordinate (F) at (1.5, 0);
    \draw[main line] (D)--(E)--(F)--cycle;
    \node[above] at (D) {$D$};
    \node[below left] at (E) {$E$};
    \node[below right] at (F) {$F$};
    \node[left]  at ($(D)!0.5!(E)$) {$3$};
    \node[below] at ($(E)!0.5!(F)$) {$4$};
    \node[right] at ($(D)!0.5!(F)$) {$5$};
    % Similar angle marks
    \draw[angle arc] ($(D)+(-100:0.4)$) arc[start angle=-100,end angle=-145,radius=0.4];
    \draw[angle arc] ($(E)+(60:0.4)$) arc[start angle=60,end angle=0,radius=0.4];
  \end{scope}

  % Similarity symbol
  \node[font=\large] at (4.7,2) {$\sim$};
  % Ratio annotation
  \node[right, align=left, font=\small] at (7.5,2.5) {$\triangle ABC \sim \triangle DEF$\\相似比 $= 2:1$};
\end{tikzpicture}
\captionof{figure}{相似三角形 $\triangle ABC \sim \triangle DEF$，相似比 $2:1$，$6:3=8:4=10:5$}
\label{fig:similar-ex1}
\end{center}


\textbf{解：}
因为 $DE \parallel BC$ ，

所以 $\triangle ADE \sim \triangle ABC$ （AA： $\angle ADE = \angle B$ ， $\angle A = \angle A$ ）。

相似比为 $\dfrac{AD}{AB} = \dfrac{3}{3 + 6} = \dfrac{3}{9} = \dfrac{1}{3}$ 。

所以 $\dfrac{DE}{BC} = \dfrac{1}{3}$ ， $BC = 3 \times DE = 3 \times 4 = 12$ 。

\textbf{例 2.} $\triangle ABC$ 中， $D$ 是 $AB$ 上的点， $E$ 是 $AC$ 上的点。 $\dfrac{AD}{AB} = \dfrac{AE}{AC} = \dfrac{2}{5}$ ， $\angle A = \angle A$ 。证明 $\triangle ADE \sim \triangle ABC$ 。

\textbf{证明：}
在 $\triangle ADE$ 和 $\triangle ABC$ 中：

$\dfrac{AD}{AB} = \dfrac{AE}{AC} = \dfrac{2}{5}$ （已知）， $\angle A = \angle A$ （公共角）。

所以 $\triangle ADE \sim \triangle ABC$ （SAS 相似）。

\textbf{例 3.} $\triangle ABC \sim \triangle DEF$ ，相似比为 $2:3$ 。 $\triangle ABC$ 的周长为 $24$ cm，面积为 $32$ cm $^2$ 。求 $\triangle DEF$ 的周长和面积。

\textbf{解：}
周长比 $=$ 相似比 $= 2:3$ 。

$\triangle DEF$ 的周长 $= 24 \times \dfrac{3}{2} = 36$ cm。

面积比 $=$ 相似比的平方 $= 4:9$ 。

$\triangle DEF$ 的面积 $= 32 \times \dfrac{9}{4} = 72$ cm $^2$ 。

\end{workedexamples}

\begin{keytakeaway}


\begin{itemize}
  \item 相似三角形：对应角相等，对应边成比例。
  \item 三种判定方法：AA、SAS 相似、SSS 相似。
  \item 周长比 $=$ 相似比；面积比 $=$ 相似比的平方。
\end{itemize}


\end{keytakeaway}

\begin{practice}


\begin{enumerate}
  \item 如图， $DE \parallel BC$ ， $AD = 4$ ， $AB = 10$ ， $BC = 15$ 。求 $DE$ 的长。
  \item $\triangle ABC \sim \triangle DEF$ ， $AB = 6$ ， $DE = 9$ ， $\triangle ABC$ 的面积为 $20$ 。求 $\triangle DEF$ 的面积。
  \item 阳光下，一根 $1.5$ m 高的木棒竖直立在地面上，影长 $2$ m。旁边一棵树的影长 $8$ m，求树的高度。
\end{enumerate}


\end{practice}



\section*{三角形的中位线}


\begin{explore}


连接三角形两边中点的线段，和第三条边有什么关系？在纸上画一个三角形，找到两边的中点连接起来，量一量这条线段和第三边——你会发现它恰好是第三边的一半，而且和第三边平行！

\end{explore}

\begin{understand}


连接三角形两边中点的线段叫做三角形的\textbf{中位线}。

\textbf{三角形中位线定理}：三角形的中位线平行于第三边，且等于第三边的一半。

\begin{center}
\begin{tikzpicture}[every node/.style={font=\small}]
  \coordinate (A) at (1, 3.5);
  \coordinate (B) at (0, 0);
  \coordinate (C) at (4.5, 0);
  \coordinate (M) at ($(A)!0.5!(B)$);
  \coordinate (N) at ($(A)!0.5!(C)$);
  \draw[main line] (A)--(B)--(C)--cycle;
  \draw[main line, draw=blue!70!black] (M)--(N);
  \node[above] at (A) {$A$};
  \node[below left]  at (B) {$B$};
  \node[below right] at (C) {$C$};
  \node[left]  at (M) {$M$};
  \node[right] at (N) {$N$};
  \fill (M) circle (2pt);
  \fill (N) circle (2pt);
  % Midpoint tick marks on AM and MB
  \draw ($(A)!0.45!(B)$)++({90+atan2(3.5,1)}:0.12)--++({-90+atan2(3.5,1)}:0.12);
  \draw ($(M)!0.5!(B)$)++({90+atan2(3.5,1)}:0.12)--++({-90+atan2(3.5,1)}:0.12);
  % Midpoint tick marks on AN and NC
  \draw ($(A)!0.45!(C)$)++({90-atan2(3.5,3.5)}:0.12)--++({-90+atan2(3.5,3.5)}:0.12);
  \draw ($(N)!0.5!(C)$)++({90-atan2(3.5,3.5)}:0.12)--++({-90+atan2(3.5,3.5)}:0.12);
  % BC label
  \node[below] at ($(B)!0.5!(C)+(0,-0.15)$) {$BC$};
  \node[blue!70!black, above] at ($(M)!0.5!(N)+(0,0.1)$) {$MN = \frac{1}{2}BC$，$MN \parallel BC$};
\end{tikzpicture}
\captionof{figure}{三角形中位线定理：$M$、$N$ 分别是 $AB$、$AC$ 中点，则 $MN \parallel BC$ 且 $MN = \frac{1}{2}BC$}
\label{fig:triangle-midline}
\end{center}


如图， $D$ 、 $E$ 分别是 $AB$ 、 $AC$ 的中点，则 $DE \parallel BC$ ，且 $DE = \dfrac{1}{2}BC$ 。

\textbf{证明}：
延长 $DE$ 到 $F$ ，使 $EF = DE$ 。连接 $CF$ 。

因为 $AE = EC$ （ $E$ 是 $AC$ 中点）， $DE = EF$ （作法）， $\angle AED = \angle CEF$ （对顶角），

所以 $\triangle ADE \cong \triangle CFE$ （SAS）。

所以 $AD = CF$ ， $\angle ADE = \angle CFE$ 。

因为 $AD = CF$ 且 $AD = DB$ （ $D$ 是 $AB$ 中点），所以 $CF = DB$ 。

因为 $\angle ADE = \angle CFE$ ，即 $\angle BDF = \angle CFD$ ，这是截线 $DF$ 所截的内错角，

所以 $BD \parallel CF$ 。

因为 $BD \parallel CF$ 且 $BD = CF$ ，所以四边形 $DBCF$ 是平行四边形（一组对边平行且相等）。

所以 $DF \parallel BC$ ，即 $DE \parallel BC$ ，且 $DF = BC$ 。

因为 $DE = \dfrac{1}{2}DF$ ，所以 $DE = \dfrac{1}{2}BC$ 。

\end{understand}

\begin{workedexamples}


\textbf{例 1.} $\triangle ABC$ 中， $D$ 、 $E$ 分别是 $AB$ 、 $AC$ 的中点， $BC = 10$ cm。求 $DE$ 的长。

\textbf{解：}
因为 $D$ 、 $E$ 分别是 $AB$ 、 $AC$ 的中点，所以 $DE$ 是 $\triangle ABC$ 的中位线。

由中位线定理， $DE = \dfrac{1}{2}BC = \dfrac{1}{2} \times 10 = 5$ cm。

\textbf{例 2.} $\triangle ABC$ 中， $D$ 、 $E$ 、 $F$ 分别是 $AB$ 、 $BC$ 、 $CA$ 的中点。 $\triangle DEF$ 的周长与 $\triangle ABC$ 的周长有什么关系？

\textbf{解：}
$DE$ 是 $\triangle ABC$ 的中位线， $DE = \dfrac{1}{2}AC$ 。

$DF$ 是 $\triangle ABC$ 的中位线， $DF = \dfrac{1}{2}BC$ 。

$EF$ 是 $\triangle ABC$ 的中位线， $EF = \dfrac{1}{2}AB$ 。

$\triangle DEF$ 的周长 $= DE + DF + EF = \dfrac{1}{2}(AC + BC + AB) = \dfrac{1}{2} \times \triangle ABC$ 的周长。

\textbf{例 3.} 四边形 $ABCD$ 中， $E$ 、 $F$ 、 $G$ 、 $H$ 分别是 $AB$ 、 $BC$ 、 $CD$ 、 $DA$ 的中点。证明 $EFGH$ 是平行四边形。

\textbf{证明：}
连接对角线 $AC$ 。

在 $\triangle ABC$ 中， $E$ 、 $F$ 分别是 $AB$ 、 $BC$ 的中点，所以 $EF \parallel AC$ ， $EF = \dfrac{1}{2}AC$ 。

在 $\triangle ACD$ 中， $G$ 、 $H$ 分别是 $CD$ 、 $DA$ 的中点，所以 $HG \parallel AC$ ， $HG = \dfrac{1}{2}AC$ 。

因为 $EF \parallel AC$ ， $HG \parallel AC$ ，所以 $EF \parallel HG$ 。

又因为 $EF = HG = \dfrac{1}{2}AC$ ，

所以四边形 $EFGH$ 是平行四边形（一组对边平行且相等）。

\end{workedexamples}

\begin{keytakeaway}


\begin{itemize}
  \item 三角形中位线平行于第三边，且等于第三边的一半。
  \item 三角形三条中位线围成的小三角形的周长是原三角形周长的一半。
  \item 中位线定理可用于证明线段平行和计算线段长度。
\end{itemize}


\end{keytakeaway}

\begin{practice}


\begin{enumerate}
  \item $\triangle ABC$ 中， $D$ 、 $E$ 分别是 $AB$ 、 $AC$ 的中点， $DE = 7$ cm。求 $BC$ 的长。
  \item 梯形 $ABCD$ 中， $AD \parallel BC$ ， $E$ 、 $F$ 分别是两腰 $AB$ 、 $CD$ 的中点。证明 $EF \parallel BC$ ，且 $EF = \dfrac{1}{2}(AD + BC)$ 。（提示：延长 $AF$ 到 $BC$ 的延长线上交于一点 $G$ ）
\end{enumerate}


\end{practice}



\begin{answer}


\begin{enumerate}
  \item $\angle B = 180° - 48° - 65° = 67°$ 。
\end{enumerate}


\begin{enumerate}
  \item 两个底角 $= \dfrac{180° - 40°}{2} = 70°$ 。
\end{enumerate}


\begin{enumerate}
  \item 由三边关系： $9 - 5 < x < 9 + 5$ ，即 $4 < x < 14$ 。第三边为整数，所以可能值为 $5, 6, 7, 8, 9, 10, 11, 12, 13$ 。
\end{enumerate}


\begin{enumerate}
  \item $4$ 不能做底边（若两腰为 $4$ ，则 $4 + 4 = 8 < 9$ ，不满足三边关系）。所以两腰为 $9$ ，底边为 $4$ ，周长 $= 9 + 9 + 4 = 22$ 。
\end{enumerate}


\begin{enumerate}
  \item 因为 $AB \parallel CD$ ，所以 $\angle ABO = \angle CDO$ （内错角相等）， $\angle BAO = \angle DCO$ （内错角相等）。又 $AB = CD$ （已知）。所以 $\triangle ABO \cong \triangle CDO$ （ASA）。
\end{enumerate}


\begin{enumerate}
  \item 在 $\triangle ADB$ 和 $\triangle BCA$ 中： $AD = BC$ （已知）， $\angle 1 = \angle 2$ （已知）， $\angle 3 = \angle 4$ （已知）。所以 $\angle 1 + \angle 3 = \angle 2 + \angle 4$ ，即 $\angle ADB = \angle BCA$ 。又 $AD = BC$ ， $AB = BA$ （公共边）。所以 $\triangle ADB \cong \triangle BCA$ （SAS）。
\end{enumerate}


\begin{enumerate}
  \item 在 $RT\triangle ABC$ 和 $RT\triangle DEF$ 中， $\angle C = \angle F = 90°$ ， $AB = DE$ （斜边相等）， $BC = EF$ （直角边相等）。所以 $\triangle ABC \cong \triangle DEF$ （HL）。
\end{enumerate}


\begin{enumerate}
  \item 情况一： $50°$ 是顶角，底角 $= \dfrac{180° - 50°}{2} = 65°$ 。其他两个角为 $65°, 65°$ 。
\end{enumerate}

   情况二： $50°$ 是底角，顶角 $= 180° - 50° - 50° = 80°$ 。其他两个角为 $50°, 80°$ 。

\begin{enumerate}
  \item 因为 $AB = AC$ ， $D$ 是 $BC$ 中点，所以 $BD = DC$ 。又 $AD = AD$ （公共边）。所以 $\triangle ABD \cong \triangle ACD$ （SSS）。所以 $\angle ADB = \angle ADC$ 。因为 $\angle ADB + \angle ADC = 180°$ ，所以 $\angle ADB = 90°$ ，即 $AD \perp BC$ 。
\end{enumerate}


\begin{enumerate}
  \item 边长 $= 18 \div 3 = 6$ cm。
\end{enumerate}


\begin{enumerate}
  \item 因为 $DE \parallel BC$ ，所以 $\triangle ADE \sim \triangle ABC$ 。 $\dfrac{DE}{BC} = \dfrac{AD}{AB} = \dfrac{4}{10} = \dfrac{2}{5}$ 。 $DE = \dfrac{2}{5} \times 15 = 6$ 。
\end{enumerate}


\begin{enumerate}
  \item 相似比 $= \dfrac{AB}{DE} = \dfrac{6}{9} = \dfrac{2}{3}$ 。面积比 $= \left(\dfrac{2}{3}\right)^2 = \dfrac{4}{9}$ 。 $\dfrac{20}{S} = \dfrac{4}{9}$ ， $S = 45$ 。
\end{enumerate}


\begin{enumerate}
  \item 木棒与其影、树与其影分别构成相似的直角三角形（阳光平行）。 $\dfrac{1.5}{2} = \dfrac{h}{8}$ ， $h = \dfrac{1.5 \times 8}{2} = 6$ m。树高 $6$ m。
\end{enumerate}


\begin{enumerate}
  \item $BC = 2 \times DE = 2 \times 7 = 14$ cm。
\end{enumerate}


\begin{enumerate}
  \item 延长 $AF$ 交 $BC$ 延长线于 $G$ 。在 $\triangle ADF$ 和 $\triangle GCF$ 中： $DF = CF$ （ $F$ 是 $CD$ 中点）， $\angle ADF = \angle GCF$ （ $AD \parallel BG$ ，内错角）， $\angle AFD = \angle GFC$ （对顶角）。所以 $\triangle ADF \cong \triangle GCF$ （ASA）， $AD = GC$ 。在 $\triangle ABG$ 中， $E$ 是 $AB$ 中点， $F$ 是 $AG$ 中点（因为 $AF = GF$ ），所以 $EF$ 是 $\triangle ABG$ 的中位线， $EF \parallel BG$ ， $EF = \dfrac{1}{2}BG = \dfrac{1}{2}(BC + CG) = \dfrac{1}{2}(BC + AD)$ 。
\end{enumerate}


\end{answer}
